\documentclass[a4paper,fontset=none]{ctexart}

\usepackage{anyfontsize}
\usepackage{amsmath,amssymb,mathrsfs,wasysym}
\usepackage{autobreak}
\allowdisplaybreaks
\usepackage[T1]{fontenc}
\usepackage{textcomp}
\usepackage{amsthm}
\usepackage[libertine,vvarbb]{newtxmath}
\usepackage[scr=rsfso,frak=euler,bb=ams]{mathalfa}
\usepackage{bm}
\usepackage{indentfirst}
\setlength{\parindent}{0em}
\usepackage{parskip}
\usepackage{geometry}
\geometry{top=3cm,bottom=3cm,left=2cm,right=2cm}
\usepackage{fontspec}
\setmainfont{Libertinus Serif}
\setsansfont{Libertinus Sans}
\setCJKmainfont{SimSun}[BoldFont=Noto Sans CJK SC Medium,ItalicFont=KaiTi]

\begin{document}

\title{\textbf{电动力学作业}}
\author{1900011604\ \ 张植竣\ \ 北京大学}
\date{2021年6月1日}
\maketitle

\section*{第七章}

\begin{itemize}

    \item[1.]
    运动电子在$\bm{n}$方向辐射场的功率为：
    \[
        \frac{\mathrm{d}P}{\mathrm{d}\varOmega} = \frac{e^2}{16\pi^2\varepsilon_0c^3} \frac{\left\|\bm{n} \times \left[(\bm{n}-\bm{\beta}) \times \dot{\bm{\beta}}\right]\right\|^2}{(1 - \bm{n}\cdot\bm{\beta})^5}
    \]
    则当$(\bm{n}-\bm{\beta})\times\dot{\bm{\beta}}=0$时，$\bm{n}$方向无辐射，即：
    \[
        \bm{n} \times \dot{\bm{\beta}} = \bm{\beta} \times \dot{\bm{\beta}}
    \]
    又因为：（其中$\alpha$为$\bm{\beta}$与$\dot{\bm{\beta}}$的夹角，$\dot{\beta}=\left\|\dot{\bm{\beta}}\right\|$）
    \[
        \left\|\bm{n} \times \dot{\bm{\beta}}\right\| = \dot{\beta}\sin\theta,\quad \left\|\bm{\beta} \times \dot{\bm{\beta}}\right\| = \beta\dot{\beta}\sin\alpha
    \]
    所以当$\sin\theta=\beta\sin\alpha=\dfrac{v}{c}\sin\alpha$时，$\bm{n}$方向无辐射。

    \item[3.(1)]
    由于$z_0\omega\ll c$，即带电粒子的速度远小于光速，辐射场可近似写为：
    \[
        \bm{E} = \frac{q}{4\pi\varepsilon_0c^2} \frac{\hat{r}\times(\hat{r}\times\dot{\bm{v}})}{r},\quad \bm{B} = \frac{1}{c}\bm{r} \times \bm{E}
    \]
    由于粒子在$t'$时刻从源点辐射的电磁波，在$t=t'+\dfrac{r}{c}$时刻才传播到场点，则有：
    \[
        z = z_0\mathrm{e}^{-i\omega t'} = z_0\mathrm{e}^{ikr-i\omega t},\quad k = \frac{\omega}{c}
    \]
    粒子的速度和加速度为：
    \[
        \bm{v} = -i\omega z_0\mathrm{e}^{ikr-i\omega t}\hat{z},\quad
        \dot{\bm{v}} = -\omega^2z_0\mathrm{e}^{ikr-i\omega t}\hat{z}
    \]
    设$\theta$为辐射方向$\hat{r}$与$\hat{z}$的夹角，由$\hat{z}=\hat{r}\cos\theta-\hat{\theta}\sin\theta$得：
    \[
        \hat{z} \times \hat{r} = \hat{\phi}\sin\theta,\quad
        \hat{r} \times (\hat{r} \times \hat{z}) = \hat{\theta}\sin\theta
    \]
    故产生的辐射场为电偶极辐射：
    \[
        \bm{E} = \frac{q}{4\pi\varepsilon_0c^2} \frac{\hat{r}\times(\hat{r}\times\dot{\bm{v}})}{r} = -\frac{q\omega^2 z_0}{4\pi\varepsilon_0c^2} \hat{\theta}\sin\theta \frac{\mathrm{e}^{ikr-i\omega t}}{r}
    \]
    \[
        \bm{B} = \frac{1}{c}\bm{r} \times \bm{E} = -\frac{q\omega^2 z_0}{4\pi\varepsilon_0c^3} \hat{\phi}\sin\theta \frac{\mathrm{e}^{ikr-i\omega t}}{r}
    \]
    辐射场的能流密度为：
    \[
        \bm{S} = \frac{1}{2\mu_0}\bm{E} \times \bm{B}^* = \frac{q^2\omega^4z_0^2}{32\pi\varepsilon_0c^3}\frac{\sin^2\theta}{r^2}\hat{r}
    \]

    \item[3.(2)]
    设地面参考系为$\varSigma$，带电粒子的静止参考系为$\varSigma'$。

    带电粒子的自场与加速度无关，在$\varSigma'$系中，粒子的自场为：
    \[
        \bm{E}' = \frac{q}{4\pi\varepsilon_0}\frac{\hat{r}'}{{r'}^2},\quad \bm{B}' = 0
    \]
    由于$z_0\omega\ll c$，即带电粒子的速度远小于光速，参考系$\varSigma'$与$\varSigma$之间的电磁场变换近似为：
    \[
        \bm{E} = \bm{E}' - \bm{v} \times \bm{B}',\quad
        \bm{B} = \bm{B}' + \frac{1}{c^2}\bm{v} \times \bm{E}'
    \]
    则在$\varSigma$系中，粒子的自场为：
    \[
        \bm{E} = \frac{q}{4\pi\varepsilon_0}\frac{\hat{r}}{r^2},\quad
        \bm{B} = \frac{q}{4\pi\varepsilon_0c^2}\frac{\bm{v}\times\hat{r}}{r^2} = -\frac{iq\omega z_0}{4\pi\varepsilon_0c^2} \hat{\phi}\sin\theta \frac{\mathrm{e}^{ikr-i\omega t}}{r^2}
    \]
    可见辐射场与$\dfrac{1}{r}$成正比，且与加速度$\dot{\bm{v}}$有关；而自场与$\dfrac{1}{r^2}$成正比，至多与速度$\bm{v}$有关。

    \item[4.]
    做圆周运动的粒子的坐标为：
    \[
        \bm{x} = R_0(\hat{x}\cos\omega t + \hat{y}\sin\omega t)
    \]
    辐射场为电偶极辐射。

    由于$\mathrm{Re}\left[(\hat{x} + i\hat{y})\mathrm{e}^{-i\omega t}\right]=\hat{x}\cos\omega t + \hat{y}\sin\omega t$，故$\bm{p}$也可以用复数表示为：
    \[
        \bm{x} = R_0(\hat{x} + i\hat{y})\mathrm{e}^{-i\omega t}
    \]
    又使用直角坐标和球坐标的转换矩阵得到：
    \[
        \bm{x} = R_0(\hat{r}\sin\theta + \theta\cos\theta + i\hat{\phi})\mathrm{e}^{i\phi-i\omega t}
    \]
    则有：
    \[
        \ddot{\bm{x}} = (-i\omega)^2\bm{x} = -\omega^2R_0(\hat{r}\sin\theta + \theta\cos\theta + i\hat{\phi})\mathrm{e}^{i\phi-i\omega t}
    \]
    由于$z_0\omega\ll c$，即带电粒子的速度远小于光速，辐射场可近似为电偶极辐射：
    \[
        \bm{E} = \frac{q}{4\pi\varepsilon_0c^2} \frac{\hat{r}\times(\hat{r}\times\dot{\bm{v}})}{r} = \frac{q\omega^2R_0}{4\pi\varepsilon_0c^2}(\hat{\theta}\cos\theta + i\hat{\phi})\frac{\mathrm{e}^{i\phi-i\omega t+ikr}}{r}
    \]
    \[
        \bm{B} = \frac{1}{c}\bm{r} \times \bm{E} = \frac{q\omega^2R_0}{4\pi\varepsilon_0c^3}(\hat{\phi}\cos\theta - i\hat{\theta})\frac{\mathrm{e}^{i\phi-i\omega t+ikr}}{r}
    \]
    辐射场的能流密度为：
    \[
        \bm{S} = \frac{1}{2\mu_0}\bm{E} \times \bm{B}^* = \frac{q^2\omega^4R_0^2}{32\pi^2\varepsilon_0c^3}(1 + \cos^2\theta)\frac{\hat{r}}{r^2}
    \]
    讨论电磁场偏振如下：
    \begin{itemize}
        \item[\bullet] $\theta=0$处：
        \[
             \bm{E} = \frac{q\omega^2R_0}{4\pi\varepsilon_0c^2}(\hat{\theta} + i\hat{\phi})\frac{\mathrm{e}^{i\phi-i\omega t+ikr}}{r}
        \]
        由于$\mathrm{Re}\left[(\hat{\theta} + i\hat{\phi})\mathrm{e}^{-i\omega t}\right]=\hat{\theta}\cos\omega t + \hat{\phi}\sin\omega t$，即$\hat{\theta}$和$\hat{\phi}$两个方向上振幅相等，但$\hat{\theta}$方向的振动比$\hat{\phi}$方向超前$\dfrac{\pi}{2}$，故在面对波传播方向的观察者看来，$\bm{E}$是左旋的圆偏振波；

        \item[\bullet] $\theta=\dfrac{\pi}{4}$处：
        \[
             \bm{E} = \frac{q\omega^2R_0}{4\pi\varepsilon_0c^2}\left(\frac{\sqrt{2}}{2}\hat{\theta} + i\hat{\phi}\right)\frac{\mathrm{e}^{i\phi-i\omega t+ikr}}{r}
        \]
        在面对波传播方向的观察者看来，$\bm{E}$是左旋的椭圆偏振波；

        \item[\bullet] $\theta=\dfrac{\pi}{2}$处：
        \[
             \bm{E} = \frac{q\omega^2R_0}{4\pi\varepsilon_0c^2}i\hat{\phi}\frac{\mathrm{e}^{i\phi-i\omega t+ikr}}{r}
        \]
        在面对波传播方向的观察者看来，$\bm{E}$是沿$\hat{\phi}$振动的线偏振波；

        \item[\bullet] $\theta=\dfrac{3\pi}{4}$处：
        \[
             \bm{E} = \frac{q\omega^2R_0}{4\pi\varepsilon_0c^2}\left(-\frac{\sqrt{2}}{2}\hat{\theta} + i\hat{\phi}\right)\frac{\mathrm{e}^{i\phi-i\omega t+ikr}}{r}
        \]
        在面对波传播方向的观察者看来，$\bm{E}$是右旋的椭圆偏振波；

        \item[\bullet] $\theta=\pi$处：
        \[
             \bm{E} = \frac{q\omega^2R_0}{4\pi\varepsilon_0c^2}(-\hat{\theta} + i\hat{\phi})\frac{\mathrm{e}^{i\phi-i\omega t+ikr}}{r}
        \]
        在面对波传播方向的观察者看来，$\bm{E}$是右旋的圆偏振波。
    \end{itemize}

    \item[5.(1)]
    设磁场方向为$\hat{z}$方向，即$\bm{B}=B\hat{z}$，则谐振子的运动方程为：
    \[
        m\ddot{\bm{r}} = -m\omega_0^2\bm{r} + q\dot{\bm{r}} \times \bm{B}
    \]
    在直角坐标系$(x,y,z)$中写为：
    \[
        m(\ddot{x}\hat{x} + \ddot{y}\hat{y} + \ddot{z}\hat{z}) = -m\omega_0^2(x\hat{x} + y\hat{y} + z\hat{z}) + qB(\dot{y}\hat{x} - \dot{x}\hat{y})
    \]
    即：
    \begin{align}
        &\ddot{x} = -\omega_0^2x + \frac{qB}{m}\dot{y} \\
        &\ddot{y} = -\omega_0^2y - \frac{qB}{m}\dot{x} \\
        &\ddot{z} = -\omega_0^2z
    \end{align}
    考虑将$\hat{r}$进行如下分解：
    \[
        \hat{r} = (x-iy)(\hat{x}+i\hat{y}) + z\hat{z}
    \]
    则由(1)式和(2)式得：
    \[
        \ddot{x} - i\ddot{y} = -\omega_0^2(x-iy) + \frac{qB}{m}(i\dot{x}+\dot{y})
    \]
    记$u=x-iy$，则有：
    \[
        \ddot{u} - i\frac{qB}{m}\dot{u} + \omega_0^2u = 0
    \]
    解得$u(t)$的通解为：
    \[
        u = C_1\mathrm{e}^{i\omega_1t} + C_2\mathrm{e}^{i\omega_2t}
    \]
    其中：
    \[
        \omega_1 = \frac{qB}{2m} + \sqrt{\left(\frac{qB}{2m}\right)^2 + \omega_0^2},\quad
        \omega_2 = \frac{qB}{2m} - \sqrt{\left(\frac{qB}{2m}\right)^2 + \omega_0^2}
    \]
    记$\omega_L=\dfrac{qB}{2m}$，在$\omega_L\ll\omega_0$的情况下（磁场对谐振子运动为微扰）：
    \[
        \omega_1 \approx \omega_L + \omega_0,\quad \omega_2 \approx \omega_L - \omega_0
    \]
    又由(3)式解得：
    \[
        z = C_3\mathrm{e}^{i\omega_0t}
    \]
    谐振子运动的通解可以写为：
    \[
        \bm{r} = C_1(\hat{x}+i\hat{y})\mathrm{e}^{-i(\omega_0-\omega_L)t} + C_2(\hat{x}+i\hat{y})\mathrm{e}^{i(\omega_L+\omega_0)t} + C_3\mathrm{e}^{i\omega_0t}
    \]
    由于$\mathrm{Re}\left[(\hat{x} + i\hat{y})\mathrm{e}^{-i\omega t}\right]=\hat{x}\cos\omega t - \hat{y}\sin\omega t$，故通解也可以写为：
    \[
        \bm{r} = \left[C_1\cos(\omega_0-\omega_L)t + C_2\cos(\omega_0+\omega_L)t\right]\hat{x} - \left[C_1\sin(\omega_0-\omega_L)t + C_2\sin(\omega_0+\omega_L)t\right]\hat{y} + (C_3\cos\omega_0t)\hat{z}
    \]

    \item[5.(2)]
    辐射场的频率和偏振讨论如下：
    \begin{itemize}
        \item[\bullet] 沿磁场方向：两个频率分别为$\omega_0-\omega_L$、$\omega_0+\omega_L$，旋转方向相反的圆偏振光；

        \item[\bullet] 垂直于磁场方向：三个频率分别为$\omega_0$、$\omega_0-\omega_L$、$\omega_0+\omega_L$的线偏振光。
    \end{itemize}

    \item[6.(1)]
    \setcounter{equation}{0}
    电子受到磁场力和辐射阻尼力，其运动方程为：
    \[
        m\ddot{\bm{r}} = e\dot{\bm{r}} \times \bm{B} + \frac{e^2}{6\pi\varepsilon_0c^3}\dddot{\bm{r}}
    \]
    记$\omega_0=\dfrac{eB}{m}$，$\lambda=\dfrac{e^2\omega_0^2}{6\pi\varepsilon_0mc^3}$，上式在直角坐标系$(x,y,z)$中写为：
    \[
        \ddot{x}\hat{x} + \ddot{y}\hat{y} + \ddot{z}\hat{z} = \omega_0(\dot{y}\hat{x} - \dot{x}\hat{y}) + \frac{\lambda}{\omega_0^2}(\dddot{x}\hat{x} + \dddot{y}\hat{y} + \dddot{z}\hat{z})
    \]
    即：
    \begin{align}
        &\ddot{x} = \omega_0\dot{y} + \frac{\lambda}{\omega_0^2}\dddot{x} \\
        &\ddot{y} = -\omega_0\dot{x} + \frac{\lambda}{\omega_0^2}\dddot{y} \\
        &\ddot{z} = \frac{\lambda}{\omega_0^2}\dddot{z}
    \end{align}
    考虑将$\hat{r}$进行如下分解：
    \[
        \hat{r} = (x-iy)(\hat{x}+i\hat{y}) + z\hat{z}
    \]
    则由(1)式和(2)式得：
    \[
        \ddot{x} - i\ddot{y} = -\omega_0^2(x-iy) + \frac{qB}{m}(i\dot{x}+\dot{y})
    \]
    记$u=x-iy$，则有：
    \begin{equation}
        \ddot{u} - i\omega_0\dot{u} = \frac{\lambda}{\omega^2}\dddot{u}
    \end{equation}
    由于：
    \[
        \frac{\lambda}{\omega_0^2} = \frac{e^2}{6\pi\varepsilon_0mc^3} = 6.2664 \times 10^{-24} \ll 1
    \]
    可以将(4)式等号右边的$\dfrac{\lambda}{\omega^2}\dddot{u}$看作微扰项略去，得到：
    \[
        \ddot{u} - i\omega_0\dot{u} = 0
    \]
    解得$u(t)$的通解为：
    \[
        u = C_1\mathrm{e}^{i\omega_0t}
    \]
    由此得$\dddot{u}=-\omega_0^2\dot{u}$，代入方程(4)得：
    \[
        \ddot{u} - (i\omega_0 - \lambda)\dot{u} = 0
    \]
    解得$u(t)$的通解为：
    \[
        u = C_1\mathrm{e}^{(i\omega_0 - \lambda)t} + C_2 = \left(C_1\mathrm{e}^{-\lambda t}\cos\omega_0t + C_2\right) + i\left(C_1\mathrm{e}^{-\lambda t}\sin\omega_0t\right)
    \]
    由于$u=x-iy$，$\omega_0\ll\lambda$，且$t=0$时的初始条件为：
    \[
        x = R_0,\quad y = 0,\quad \dot{x} = 0,\quad \dot{y} = v_0
    \]
    得$x$和$y$坐标的运动轨迹：
    \begin{align}
        &x = v_0\omega_0\mathrm{e}^{-\lambda t}\cos\omega_0t + \left(R_0 - v_0\omega_0\right) \\
        &y = v_0\omega_0\mathrm{e}^{-\lambda t}\sin\omega_0t
    \end{align}
    又由(3)式解得：
    \[
        z = C_1\mathrm{e}^{\tfrac{\omega_0^2}{\lambda}t} + C_2x + C_3
    \]
    而$t=0$时的初始条件为：
    \[
        z = 0,\quad \dot{z} = 0
    \]
    则$z$坐标的运动轨迹为：
    \begin{equation}
        z = 0
    \end{equation}
    综合(5)、(6)、(7)三式可得：电子运动轨迹是一条$Oxy$平面上的螺旋线。

    \item[6.(2)]
    由(5)、(6)、(7)三式可得电子的加速度分量：
    \begin{align}
        &\ddot{x} = v_0\omega_0\left[(\lambda^2-1)\mathrm{e}^{-\lambda t}\cos\omega_0t + 2\lambda\mathrm{e}^{-\lambda t}\sin\omega_0t\right] \\
        &\ddot{y} = v_0\omega_0\left[(\lambda^2-1)\mathrm{e}^{-\lambda t}\sin\omega_0t - 2\lambda\mathrm{e}^{-\lambda t}\cos\omega_0t\right] \\
        &\ddot{z} = 0
    \end{align}
    由于$\lambda\ll\omega_0$，可将$\lambda$项和$\lambda^2$项均略去，得到：
    \[
        \ddot{x} = -v_0\omega_0\mathrm{e}^{-\lambda t}\cos\omega_0t,\quad \ddot{y} = -v_0\omega_0\mathrm{e}^{-\lambda t}\sin\omega_0t,\quad \ddot{z} = 0
    \]
    则电子的加速度大小为：
    \[
        \dot{\bm{v}} = \sqrt{\ddot{x}^2 + \ddot{y}^2 + \ddot{z}^2} = v_0^2\omega_0^2\mathrm{e}^{-2\lambda t}
    \]
    在非相对论条件下，由Larmor公式得电子单位时间的辐射能量（辐射功率）为：
    \[
        P = \frac{e^2\dot{\bm{v}}^2}{6\pi\varepsilon_0c^3} = \frac{e^2v_0^2\omega_0^2}{6\pi\varepsilon_0c^3}\mathrm{e}^{-2\lambda t}
    \]

    \item[7.(1)]
    \setcounter{equation}{0}
    由$\gamma=\dfrac{1}{\sqrt{1-\bm{\beta}\cdot\bm{\beta}}}$得：
    \[
        \frac{\mathrm{d}\gamma}{\mathrm{d}t} = \frac{\mathrm{d}\gamma}{\mathrm{d}\bm{\beta}} \cdot \frac{\mathrm{d}\bm{\beta}}{\mathrm{d}t} = (\bm{\beta}\gamma^3) \cdot \dot{\bm{\beta}} = (\bm{\beta} \cdot \dot{\bm{\beta}})\gamma^3
    \]
    利用乘积的求导法则可得：（其中$\mathrm{ret}$表示时刻$t'=t-\dfrac{r}{c}$时的值）
    \[
        \bm{F} = \frac{\mathrm{d}}{\mathrm{d}t}(\gamma m\bm{v}) = \frac{\mathrm{d}}{\mathrm{d}t}(\gamma mc\bm{\beta}) = \gamma mc\left[\dot{\bm{\beta}} + \gamma^2(\bm{\beta} \cdot \dot{\bm{\beta}})\bm{\beta}\right]_{\mathrm{ret}}
    \]
    于是有：（注意$1+\gamma^2\beta^2=1+\dfrac{\beta^2}{1-\beta^2}=\dfrac{1}{1-\beta^2}=\gamma^2$以及$\bm{\beta}\times\bm{\beta}=0$）
    \[
        \bm{\beta} \cdot \bm{F} = \gamma mc\bm{\beta} \cdot \left[\dot{\bm{\beta}} + \gamma^2(\bm{\beta} \cdot \dot{\bm{\beta}})\bm{\beta}\right] = \gamma mc (\bm{\beta} \cdot \dot{\bm{\beta}}) (1+\gamma^2\beta^2) = \gamma^3mc(\bm{\beta} \cdot \dot{\bm{\beta}})
    \]
    \[
        \bm{\beta} \times \bm{F} = \gamma mc\bm{\beta} \times \left[\dot{\bm{\beta}} + \gamma^2(\bm{\beta} \cdot \dot{\bm{\beta}})\bm{\beta}\right] = \gamma mc(\bm{\beta} \times \dot{\bm{\beta}})
    \]
    \[
        \hat{r} \times \bm{F} = \gamma mc\hat{r} \times \left[\dot{\bm{\beta}} + \gamma^2(\bm{\beta} \cdot \dot{\bm{\beta}})\bm{\beta}\right] = \gamma mc(\hat{r} \times \dot{\bm{\beta}}) + \gamma^3mc(\bm{\beta} \cdot \dot{\bm{\beta}})\hat{r} \times \bm{\beta}
    \]
    则有：
    \[
        (\hat{r}-\bm{\beta}) \times \bm{F} - \bm{\beta} \cdot \bm{F}(\hat{r} \times \bm{\beta}) = \gamma mc(\hat{r}-\bm{\beta}) \times \dot{\bm{\beta}}
    \]
    即：
    \begin{equation}
        (\hat{r}-\bm{\beta}) \times \dot{\bm{v}} = \frac{1}{\gamma m}\left[(\hat{r}-\bm{\beta}) \times \bm{F} - \bm{\beta} \cdot \bm{F}(\hat{r} \times \bm{\beta})\right]
    \end{equation}
    故带电粒子的辐射场公式：
    \[
        \bm{E} = \frac{q}{4\pi\varepsilon_0c^2r} \frac{(\hat{r}-\bm{\beta}) \times \dot{\bm{v}}}{(1 - \hat{r} \cdot \bm{\beta})^3}
    \]
    用作用力可以表示为：
    \[
        \bm{E} = \frac{q}{4\pi\varepsilon_0c^2r}\left\{\frac{\delta^3}{\gamma m}\hat{r} \times \left[(\hat{r}-\bm{\beta}) \times \bm{F} - \bm{\beta} \cdot \bm{F}(\hat{r} \times \bm{\beta})\right]\right\}_{\mathrm{ret}}
    \]
    其中$\delta=\dfrac{1}{1-\hat{r}\cdot\bm{\beta}}$。

    \item[7.(2)]
    利用公式$(\bm{A}\times\bm{B})^2=A^2B^2-(\bm{A}\cdot\bm{B})^2$得：
    \begin{equation}
        \begin{aligned}
            \left[(\hat{r}-\bm{\beta}) \times \bm{F}\right]^2
            &= (\hat{r}-\bm{\beta})^2 F^2 - \left[(\hat{r}-\bm{\beta}) \cdot \bm{F}\right]^2 \\
            &= (\hat{r}-\bm{\beta})^2 F^2 - (\hat{r} \cdot \bm{F} - \bm{\beta} \cdot \bm{F})^2 \\
            &= (1 - 2\hat{r}\cdot\bm{\beta} + \beta^2) F^2 - (\hat{r} \cdot \bm{F})^2 - (\bm{\beta} \cdot \bm{F})^2 + 2(\hat{r} \cdot \bm{F})(\bm{\beta} \cdot \bm{F}) \\
            &= F^2 - 2(\hat{r}\cdot\bm{\beta})F^2 + \beta^2F^2 - (\hat{r} \cdot \bm{F})^2 - (\bm{\beta} \cdot \bm{F})^2 + 2(\hat{r} \cdot \bm{F})(\bm{\beta} \cdot \bm{F})
        \end{aligned}
    \end{equation}
    以及：
    \begin{equation}
        \begin{aligned}
            \left[\bm{F} \cdot (\hat{r} \times \bm{\beta})\right]^2
            &= F^2 (\hat{r} \times \bm{\beta})^2 - \left[\bm{F} \times (\hat{r} \times \bm{\beta})\right]^2 \\
            &= F^2 \left[r^2\beta^2 - (\hat{r}\cdot\bm{\beta})^2\right] - \left[(\bm{F}\cdot\bm{\beta})\hat{r} - (\bm{F}\cdot\hat{r})\bm{\beta}\right]^2 \\
            &= F^2\beta^2 - F^2 (\hat{r} \cdot \bm{\beta})^2 - (\bm{F} \cdot \bm{\beta})^2 - \beta^2 (\bm{F} \cdot \hat{r})^2 + 2(\bm{F} \cdot \bm{\beta})(\bm{F} \cdot \hat{r})(\hat{r} \cdot \bm{\beta}) \\
            &= \beta^2F^2 - (\hat{r} \cdot \bm{\beta})^2 F^2 - (\bm{\beta} \cdot \bm{F})^2 - \beta^2 (\hat{r} \cdot \bm{F})^2 + 2(\hat{r} \cdot \bm{\beta})(\hat{r} \cdot \bm{F})(\bm{\beta} \cdot \bm{F})
        \end{aligned}
    \end{equation}
    则(2)、(3)两式即为所求。

    \item[7.(3)]
    由(1)式得：
    \[
        \left[\hat{r} \times (\hat{r}-\bm{\beta}) \times \dot{\bm{v}}\right]^2 = \left\{\frac{1}{\gamma m} \hat{r} \times \left[(\hat{r}-\bm{\beta}) \times \bm{F} - (\bm{\beta} \cdot \bm{F})\hat{r} \times \bm{\beta}\right]\right\}^2
    \]
    即：
    \begin{equation}
        \gamma^2m^2 \left[\hat{r} \times (\hat{r}-\bm{\beta}) \times \dot{\bm{v}}\right]^2 = \left[\hat{r} \times (\hat{r}-\bm{\beta}) \times \bm{F} - (\bm{\beta}\cdot\bm{F}) \hat{r} \times (\hat{r} \times \bm{\beta})\right]^2
    \end{equation}
    又由$\bm{a}\cdot\bm{b}\times\bm{c}=\bm{c}\cdot\bm{a}\times\bm{b}$得：
    \begin{align*}
        \left[\hat{r} \times (\hat{r}-\bm{\beta}) \times \bm{F}\right]^2
        &= \hat{r}^2 \left[(\hat{r}-\bm{\beta}) \times \bm{F}\right]^2 - \left[\hat{r} \cdot (\hat{r}-\bm{\beta}) \times \bm{F}\right]^2 \\
        &= \left[(\hat{r}-\bm{\beta}) \times \bm{F}\right]^2 - \left[\bm{F} \cdot \hat{r} \times (\hat{r}-\bm{\beta})\right]^2 \\
        &= \left[(\hat{r}-\bm{\beta}) \times \bm{F}\right]^2 - \left[\bm{F} \cdot (\hat{r}\times\hat{r} - \hat{r}\times\bm{\beta})\right]^2 \\
        &= \left[(\hat{r}-\bm{\beta}) \times \bm{F}\right]^2 - \left[\bm{F} \cdot (\hat{r} \times \bm{\beta})\right]^2
    \end{align*}
    再代入(2)、(3)两式得：
    \begin{equation}
        \begin{aligned}
            &\quad \left[\hat{r} \times (\hat{r}-\bm{\beta}) \times \bm{F}\right]^2 \\
            &= \left[1 - 2(\hat{r} \cdot \bm{\beta}) + (\hat{r} \cdot \bm{\beta})^2\right] F^2 + (1 - \beta^2) (\hat{r} \cdot \bm{F})^2 + 2\left[1 - (\hat{r} \cdot \bm{\beta})\right] (\hat{r} \cdot \bm{F}) (\bm{\beta} \cdot \bm{F}) \\
            &= \frac{F^2}{\delta^2} - \frac{1}{\gamma^2}(\hat{r} \cdot \bm{F})^2 + \frac{2}{\delta}(\hat{r} \cdot \bm{F})(\bm{\beta} \cdot \bm{F})
        \end{aligned}
    \end{equation}
    其中$\delta=\dfrac{1}{1-\hat{r}\cdot\bm{\beta}}$，$\gamma=\dfrac{1}{\sqrt{1-\beta^2}}$。

    又有：
    \begin{equation}
        \begin{aligned}
            &\quad \left[(\bm{\beta} \cdot \bm{F})\hat{r} \times (\hat{r} \times \bm{\beta})\right]^2 \\
            &= (\bm{\beta} \cdot \bm{F})^2 \left[(\hat{r} \cdot \bm{\beta})\hat{r} - (\hat{r} \cdot \hat{r})\bm{\beta}\right]^2 \\
            &= (\bm{\beta} \cdot \bm{F})^2 \left[(\hat{r} \cdot \bm{\beta})\hat{r} - \bm{\beta}\right]^2 \\
            &= (\bm{\beta} \cdot \bm{F})^2 \left[(\hat{r} \cdot \bm{\beta})^2\hat{r}^2 + \beta^2 - 2(\hat{r} \cdot \bm{\beta}) (\hat{r} \cdot \bm{\beta})\right]^2 \\
            &= (\bm{\beta} \cdot \bm{F})^2 \left[\beta^2 - (\hat{r} \cdot \bm{\beta})^2\right] \\
            &= \beta^2 (\bm{\beta} \cdot \bm{F})^2 - (\hat{r} \cdot \bm{\beta})^2 (\bm{\beta} \cdot \bm{F})^2
        \end{aligned}
    \end{equation}
    以及：
    \begin{align*}
        &\quad \left[\hat{r} \times (\hat{r}-\bm{\beta}) \times \bm{F}\right] \cdot \left[(\bm{\beta} \cdot \bm{F})\hat{r} \times (\hat{r} \times \bm{\beta})\right] \\
        &= (\bm{\beta} \cdot \bm{F}) \left[\hat{r} \times (\hat{r}-\bm{\beta}) \times \bm{F}\right] \cdot \left[\hat{r} \times (\hat{r} \times \bm{\beta})\right] \\
        &= (\bm{\beta} \cdot \bm{F}) \left\{(\hat{r} \cdot \bm{F}) (\hat{r} - \hat{\beta}) - \left[\hat{r} \cdot (\hat{r} - \bm{\beta})\right] \bm{F}\right\} \cdot \left[(\hat{r} - \bm{\beta}) \hat{r} - \bm{\beta}\right] \\
        &= (\bm{\beta} \cdot \bm{F}) \left\{(\hat{r} \cdot \bm{F}) (\hat{r} - \hat{\beta}) \cdot \hat{r} - \left[\hat{r} \cdot (\hat{r} - \bm{\beta})\right] \bm{F} \cdot \hat{r} - (\hat{r} \cdot \bm{\beta} - \beta^2) \hat{r} \cdot \bm{F} + \left[\hat{r} \cdot (\hat{r} - \bm{\beta})\right] \bm{\beta} \cdot \bm{F}\right\} \\
        &= (\bm{\beta} \cdot \bm{F}) \left\{\left[\hat{r} \cdot (\hat{r} - \bm{\beta})\right] \bm{\beta} \cdot \bm{F} - (\hat{r} \cdot \bm{\beta} - \beta^2) \hat{r} \cdot \bm{F}\right\} \\
        &= (\bm{\beta} \cdot \bm{F}) \left\{(1 - \hat{r} \cdot \bm{\beta}) \bm{\beta} \cdot \bm{F} - (\hat{r} \cdot \bm{\beta} - \beta^2) \hat{r} \cdot \bm{F}\right\} \\
        &= \frac{1}{\delta} (\bm{\beta} \cdot \bm{F})^2 - (\hat{r} \cdot \bm{\beta} - \beta^2) (\hat{r} \cdot \bm{F}) (\bm{\beta} \cdot \bm{F})
    \end{align*}
    注意：
    \[
        \hat{r} \cdot \bm{\beta} - \beta^2 = (1 - \beta^2) - (1 - \hat{r} \cdot \bm{\beta}) = \frac{1}{\gamma^2} - \frac{1}{\delta}
    \]
    则有：
    \begin{equation}
        \quad \left[\hat{r} \times (\hat{r}-\bm{\beta}) \times \bm{F}\right] \cdot \left[(\bm{\beta} \cdot \bm{F})\hat{r} \times (\hat{r} \times \bm{\beta})\right] = \frac{1}{\delta} (\bm{\beta} \cdot \bm{F})^2 - \frac{1}{\gamma^2} (\hat{r} \cdot \bm{F}) (\bm{\beta} \cdot \bm{F}) + \frac{1}{\delta} (\hat{r} \cdot \bm{F}) (\bm{\beta} \cdot \bm{F})
    \end{equation}
    将(5)、(6)、(7)式代入(4)式得到：
    \begin{align*}
        &\quad \gamma^2m^2 \left[\hat{r} \times (\hat{r}-\bm{\beta}) \times \dot{\bm{v}}\right]^2 \\
        &= \frac{F^2}{\delta^2} - \frac{1}{\gamma^2}(\hat{r} \cdot \bm{F})^2 + \beta^2 (\bm{\beta} \cdot \bm{F})^2 - (\hat{r} \cdot \bm{\beta})^2 (\bm{\beta} \cdot \bm{F})^2 - \frac{2}{\delta} (\bm{\beta} \cdot \bm{F})^2 + \frac{2}{\gamma^2} (\hat{r} \cdot \bm{F}) (\bm{\beta} \cdot \bm{F}) \\
        &= \frac{F^2}{\delta^2} - \left[(\hat{r} \cdot \bm{\beta})^2 - \beta^2 + \frac{2}{\delta}\right] (\bm{\beta} \cdot \bm{F})^2 - \frac{1}{\gamma^2}(\hat{r} \cdot \bm{F})^2 + \frac{2}{\gamma^2} (\hat{r} \cdot \bm{F}) (\bm{\beta} \cdot \bm{F})
    \end{align*}
    因为：
    \begin{align*}
        (\hat{r} \cdot \bm{\beta})^2 - \beta^2 + \frac{2}{\delta}
        &= (\hat{r} \cdot \bm{\beta})^2 - \beta^2 + 2 - 2(\hat{r} \cdot \bm{\beta}) \\
        &= 1 - 2(\hat{r} \cdot \bm{\beta}) + (\hat{r} \cdot \bm{\beta})^2 + 1 - \beta^2 \\
        &= \left[1 - (\hat{r} \cdot \bm{\beta})\right]^2 + \left(1 - \beta^2\right) \\
        &= \frac{1}{\delta^2} + \frac{1}{\gamma^2}
    \end{align*}
    则有：
    \begin{align*}
        \gamma^2m^2 \left[\hat{r} \times (\hat{r}-\bm{\beta}) \times \dot{\bm{v}}\right]^2
        &= \frac{F^2}{\delta^2} - \left(\frac{1}{\delta^2} + \frac{1}{\gamma^2}\right) (\bm{\beta} \cdot \bm{F})^2 - \frac{1}{\gamma^2}(\hat{r} \cdot \bm{F})^2 + \frac{2}{\gamma^2} (\hat{r} \cdot \bm{F}) (\bm{\beta} \cdot \bm{F}) \\
        &= \frac{F^2}{\delta^2} - \frac{1}{\delta^2} (\bm{\beta} \cdot \bm{F})^2 - \frac{1}{\gamma^2} \left[(\hat{r} \cdot \bm{F})^2 + (\bm{\beta} \cdot \bm{F})^2 - 2 (\hat{r} \cdot \bm{F}) (\bm{\beta} \cdot \bm{F})\right] \\
        &= \frac{\bm{F}^2}{\delta^2} - \frac{1}{\delta^2}(\bm{\beta} \cdot \bm{F})^2 - \frac{1}{\gamma^2}(\hat{r} \cdot \bm{F} - \bm{\beta} \cdot \bm{F})^2
    \end{align*}
    故带电粒子辐射功率的角分布公式：
    \[
        \frac{\mathrm{d}P}{\mathrm{d}\varOmega} = \frac{q^2}{16\pi^2\varepsilon_0c^3} \frac{\left[\hat{r} \times (\hat{r}-\bm{\beta}) \times \dot{\bm{v}}\right]^2}{(1 - \bm{\beta} \cdot \hat{r})^5} = \frac{q^2}{16\pi^2\varepsilon_0c^3} \cdot \delta^5\left[\hat{r} \times (\hat{r}-\bm{\beta}) \times \dot{\bm{v}}\right]^2
    \]
    用作用力可以表示为：
    \[
        \frac{\mathrm{d}P}{\mathrm{d}\varOmega} = \frac{q^2}{16\pi^2\varepsilon_0c^3} \cdot \frac{\delta^3}{\gamma^2m^2}\left[\bm{F}^2 - (\bm{\beta} \cdot \bm{F})^2 - \frac{\delta^2}{\gamma^2}(\hat{r} \cdot \bm{F} - \bm{\beta} \cdot \bm{F})^2\right]
    \]
    其中$\delta=\dfrac{1}{1-\hat{r}\cdot\bm{\beta}}$。

    \item[9.(1)]
    \setcounter{equation}{0}
    粒子受到的Lorentz力提供加速度：
    \begin{equation}
        q\bm{v} \times \bm{B} = \frac{\mathrm{d}\bm{p}}{\mathrm{d}t}
    \end{equation}
    即：（此时$\bm{v}$大小不变，则$\gamma$为常数）
    \begin{equation}
        q\bm{\beta} \times \bm{B} = \gamma m\dot{\bm{\beta}}
    \end{equation}
    相对论性带电粒子辐射总功率用Li\'{e}nard公式描述：
    \[
        P = \frac{q^2}{6\pi\varepsilon_0c}\gamma^6\left[\dot{\beta}^2 - (\bm{\beta} \times \dot{\bm{\beta}})^2\right]
    \]
    此时粒子做圆周运动，$\bm{\beta}\perp\dot{\bm{\beta}}$，则有：
    \[
        \dot{\beta}^2 - (\bm{\beta} \times \dot{\bm{\beta}})^2 = \dot{\beta}^2 - (\beta\dot{\beta})^2 = (1 - \beta^2)\dot{\beta}^2 = \frac{\dot{\beta}^2}{\gamma^2}
    \]
    于是得到带电粒子的辐射功率：
    \[
        P = \frac{q^2\gamma^2}{6\pi\varepsilon_0m^2c^3}\left(\frac{\mathrm{d}\bm{p}}{\mathrm{d}t}\right)^2 = \frac{q^2}{6\pi\varepsilon_0c}\gamma^4\dot{\beta}^2
    \]
    代入(1)或(2)式即得：（注意$\beta^2=\dfrac{\gamma^2-1}{\gamma^2}$）
    \[
        P = \frac{q^4B^2(\gamma^2-1)}{6\pi\varepsilon_0m^2c}
    \]

    \item[9.(2)]
    粒子的能量为$E=\gamma mc^2$，而其能量减少来自于辐射：
    \[
        \frac{\mathrm{d}E}{\mathrm{d}t} = -\frac{q^4B^2(\gamma^2-1)}{6\pi\varepsilon_0m^2c} = -\frac{q^4B^2}{6\pi\varepsilon_0m^2c} \left(\frac{E^2}{m^2c^4}-1\right)
    \]
    即：
    \[
        \frac{\mathrm{d}E}{(E^2/m^2c^4)-1} = -\frac{q^4B^2}{6\pi\varepsilon_0m^2c}\,\mathrm{d}t
    \]
    由积分公式：
    \[
        \int \frac{\mathrm{d}y}{y^2-1} = \int \frac{1}{2}\left(\frac{1}{y+1} + \frac{1}{y-1}\right)\,\mathrm{d}y = \frac{1}{2}\ln\left|\frac{1+y}{1-y}\right|
    \]
    以及初始条件$t=t_0$时$E_0=\gamma_0mc^2$，解得：
    \[
        E(t) = mc^2 \cdot \frac{\gamma_0 + 1 + (\gamma_0 - 1)\mathrm{e}^{-\eta(t-t_0)}}{\gamma_0 + 1 - (\gamma_0 - 1)\mathrm{e}^{-\eta(t-t_0)}},\quad \eta = \frac{q^4B^2}{3\pi\varepsilon_0m^3c^3}
    \]

    \item[9.(3)]
    在非相对论条件下，电子辐射的功率退化为Larmor公式：
    \[
        P = \frac{e^2\dot{\bm{v}}^2}{6\pi\varepsilon_0c^3}
    \]
    而Lorentz力的公式为：
    \[
        q\bm{v} \times \bm{B} = m\dot{\bm{v}}
    \]
    由题意$\bm{v}\perp\bm{B}$，于是有：
    \[
        P = \frac{e^2}{6\pi\varepsilon_0c^3}\left(\frac{qvB}{m}\right)^2 = \frac{q^4v^2B^2}{6\pi\varepsilon_0m^2c^3}
    \]
    在非相对论条件下，动能$T=\dfrac{1}{2}mv^2$，则：
    \[
        \frac{\mathrm{d}T}{\mathrm{d}t} = -\frac{q^4v^2B^2}{6\pi\varepsilon_0m^2c^3} = -\frac{q^4B^2T}{3\pi\varepsilon_0m^3c^3}
    \]
    即：
    \[
        \frac{\mathrm{d}T}{T} = -\frac{q^4v^2B^2}{6\pi\varepsilon_0m^2c^3} = -\frac{q^4B^2}{3\pi\varepsilon_0m^3c^3}\,\mathrm{d}t
    \]
    由初始条件$t=t_0$时$T=T_0$，解得：
    \[
        T = T_0\mathrm{e}^{-\eta(t-t_0)},\quad \eta = \frac{q^4B^2}{3\pi\varepsilon_0m^3c^3}
    \]

\end{itemize}

\subsection*{补充题}

\begin{itemize}

    \item[2.]
    粒子受到的Lorentz力提供圆周运动的向心力：
    \[
        e\bm{v} \times \bm{B} = -\frac{mv^2}{r}\hat{r}
    \]
    在相对论条件下：
    \[
        m = \gamma m_0,\quad \bm{p} = \gamma m_0\bm{v}
    \]
    于是有：
    \[
        \frac{epB}{\gamma m_0} = \frac{\gamma m_0}{r}\left(\frac{p}{\gamma m_0}\right)^2 = \frac{p^2}{\gamma m_0r}
    \]
    解得轨道半径为：
    \[
        r = \frac{p}{eB}
    \]
    说明非相对论条件下的轨道半径公式$r=\dfrac{p}{eB}$可以推广到相对论条件，但此处$p=mv=\gamma m_0v$为相对论条件下的动量。

    辐射总功率为：
    \[
        P = \frac{q^2\gamma^2}{6\pi\varepsilon_0m^2c^3}\left(\frac{\mathrm{d}\bm{p}}{\mathrm{d}t}\right)^2 = \frac{q^2\gamma^2}{6\pi\varepsilon_0m^2c^3} \left(\frac{epB}{\gamma m_0}\right)^2 = \frac{e^4p^2B^2}{6\pi\varepsilon_0m^4c^3}
    \]

    \item[3.]
    电子的静能为$m_0c^2=0.511\,\mathrm{MeV}$，加速到$10\,\mathrm{MeV}$时：
    \[
        \gamma = \frac{E}{m_0c^2} = \frac{10\,\mathrm{MeV}}{0.511\,\mathrm{MeV}} = 19.5695
    \]
    \[
        \beta = \sqrt{\frac{\gamma^2-1}{\gamma^2}} = 0.9987
    \]
    在非相对论条件下，电子的辐射功率为：
    \[
        P_{\mathrm{nr}} = \frac{\dot{\beta}^2}{6\pi\varepsilon_0c}
    \]
    当$\dot{\bm{v}}\varparallel\bm{v}$时：
    \[
        P = \frac{\gamma^6\dot{\beta}^2}{6\pi\varepsilon_0c} = \gamma^6P_{\mathrm{nr}}
    \]
    则有：
    \[
        \frac{P}{P_{\mathrm{nr}}} = \gamma^6 = 5.6167 \times 10^7
    \]
    当$\dot{\bm{v}}\perp\bm{v}$时：
    \[
        P = \frac{\gamma^4\dot{\beta}^2}{6\pi\varepsilon_0c} = \gamma^4P_{\mathrm{nr}}
    \]
    则有：
    \[
        \frac{P}{P_{\mathrm{nr}}} = \gamma^4 = 1.4666 \times 10^5
    \]

\end{itemize}

\end{document}